\scanpage{246}
\[
\tau_{N',N}(M)\hookrightarrow M
\]
Soit $M\in\operatorname{Ob}\mathcal M_N$, $\unsure{P}\in\operatorname{Ob}\mathcal M_{N'}$,
\[
\diagram{\begin{tikzcd}[column sep=large]
\operatorname{Hom}_{\mathcal M_N}(i_{N,N'}(P),M)\arrow[d,equal]&\operatorname{Hom}_{\mathcal M_{N'}}(P,\tau_{N',N}(M))\arrow[l,"\sim"']\arrow[d,equal]\\
\operatorname{Hom}_{\mathcal M}(P,M)&\operatorname{Hom}(P,\tau_{N',N}(M))\arrow[l,"\sim"']
\end{tikzcd}}
\]
\uline{Adjoint à gauche} $q_{N',N}$ \uline{de} $i_{N,N'}$
\[
\diagram{\begin{tikzcd}[column sep=large]
\operatorname{Hom}_{\mathcal M_N}(M,i_{N,N'}(\unsure{P}))=\arrow[d,equal]&\\
\operatorname{Hom}_{\mathcal M}(M,P)\arrow[r,equal]&\operatorname{Hom}_{\mathcal M}(M/\substack{\text{ss.-module}\\\text{engendré par}\\\text{les }M_i,\ i<N'},P
\end{tikzcd}}
\]
\[
\diagram{\begin{array}{cccccccccc}
0&0&M_N&M_{N+1}&M_{N+2}&\cdots&M_{N'-1}&M_{N'}&M_{N'+1}&\cdots\\
0&0&0&&0&\cdots&0&\unsure{P}_{N'}&\unsure{P}_{N'+1}&\cdots
\end{array}}
\]
\[
q_{N',N}(M)=M/\substack{\text{ss.-module de }M\text{ engendré}\\\text{par les }M_i,\ i<N'}
\]
\uline{NB} $q_{N',N}$ et $\tau_{N',N}$ sont des ``foncteurs de localisation'' (car \unsure{adjoints} d'un foncteur pl. fid.);
\deleted{et} $q_{N',N}$ est exact à droite (en fait, commute aux $\varinjlim$ \unsure{quelc.}) mais en général pas à gauche, sans doute --- on ne s'attend pas qu'il ait encore un adj. à g. Par contre $\tau_{N',N}$ commute aux $\varprojlim$ \uline{et} aux $\varinjlim$ (à fortiori il est exact). On s'attend à en trouver un foncteur adjoint à droite.
\margin{$q_{N',N},\ i_{N,N'},\ \tau_{N',N}$}
\margin{245}
